\chapter{12.10-2 等腰三角形判定}

\section*{一、选择题（共 8 小题）}

\questionlist[8]
\item
  \begin{questionwithimage}{12.10-1-8.png}
  \item 在 $4 \times 4$ 的方格表中给出如图所示的 8 个点，任选三个作为三角形的顶点，共可构成（\quad）个等腰三角形。
  \optionlist
  \item 5
  \item 6
  \item 7
  \item 8
  \end{enumerate}
  \end{questionwithimage}


\end{enumerate}

\section*{二、填空题（共 6 小题）}

\blanklist[9]

\item 已知 $a, b, c$ 为三个正整数，如果 $a + b + c = 12$，那么以 $a, b, c$ 为边能组成的三角形是：①等腰三角形；②等边三角形；③直角三角形；④钝角三角形。以上符合条件的正确结论是\_\_\_\_\_\_。（只填序号）

\item 如图，$AD$ 是 $\triangle ABC$ 的边 $BC$ 上的高，现给出下列条件：① $\angle BAD = \angle ACD$；② $\angle BAD = \angle CAD$；③ $BD = CD$；④ $AB + BD = AC + CD$，若添加这些条件中的某一个就能推出 $\triangle ABC$ 是等腰三角形，这个条件可以是\_\_\_\_\_\_。（只填序号）

\rightquestionimage[width=0.35\textwidth]{0.35\textwidth}{12.10-11.png}

\item 如图，$AD$ 是直角三角形 $\triangle ABC$ 斜边上的中线，把 $ADC$ 沿 $AD$ 对折，点 $C$ 落在点 $C'$ 处，连接 $CC'$，则图中共有等腰三角形\_\_\_\_\_\_个。

\rightquestionimage[width=0.35\textwidth]{0.35\textwidth}{12.10-12.png}

\item 如图，组成虚线网格的每个小正三角形的边长都为 1，若有格点（小正三角形的顶点）$C$，使 $\triangle ABC$ 为等腰三角形，那么这样的 $C$ 点共有\_\_\_\_\_\_个。

\rightquestionimage[width=0.35\textwidth]{0.35\textwidth}{12.10-13.png}

\item 以原点 $O(0, 0)$，$A(1, 2)$ 为顶点的 $\triangle AOB$ 是等腰三角形，且点 $B$ 在坐标轴上，满足条件的点 $B$ 有\_\_\_\_\_\_个，其中横纵坐标均为整数的坐标为\_\_\_\_\_\_。

\item 如图，点 $A$ 的坐标为 $(0, 8)$，点 $B$ 的坐标为 $(6, 0)$，在 $x$ 轴上确定一点 $P$，使 $\triangle PAB$ 为一个等腰三角形，则 $R$ 点的坐标可以是\_\_\_\_\_\_。

\rightquestionimage[width=0.35\textwidth]{0.35\textwidth}{12.10-14.png}

\end{enumerate}

\section*{三、解答题（共 4 小题）}

\solutionlist[15]

\item 如图：已知 $AB = AE$，$BC = ED$，$\angle B = \angle E$，$AF \perp CD$，$F$ 为垂足，求证：① $AC = AD$；② $CF = DF$。

\rightquestionimage[width=0.35\textwidth]{0.35\textwidth}{12.10-15.png}

\vspace{2cm}

\item 已知，如图，在 $\triangle ABC$ 中，$AD \perp BC$，垂足为点 $D$，$BE \perp AC$，垂足为点 $E$，$M$ 为 $AB$ 边的中点，连接 $ME$、$MD$、$ED$。

(1) 求证：$\triangle MED$ 为等腰三角形；

(2) 求证：$\angle EMD = 2\angle DAC$。

\rightquestionimage[width=0.35\textwidth]{0.35\textwidth}{12.10-16.png}

\vspace{5cm}

\item 已知：一张直角三角形纸片如图1放置在平面直角坐标系中，一条直角边 $OA$ 落在 $x$ 轴正半轴上，另一条直角边 $OB$ 落在 $y$ 轴正半轴上，且 $OA = 8$，$OB = 6$。现再找一个与 $Rt\triangle ABO$ 有一条公共边且不重叠的三角形，使它们拼在一起后能构成一个大的等腰三角形。例如：如图2，$\triangle CBO$ 与 $\triangle ABO$ 拼成等腰 $\triangle ABC$，则点 $C$ 坐标为 $(-2, 0)$。请直接写出除图2情况外，其他所有的所拼成的等腰三角形中除 $A$、$B$、$O$ 三点外另一顶点 $P$ 的坐标。

\rightquestionimage[width=0.50\textwidth]{0.50\textwidth}{12.10-17.png}

\vspace{5cm}

\end{enumerate}
